\section{Discussion}
\label{Discussion}
This section interprets the results, connects them to the underlying mechanisms, and discusses limitations of the study.

\subsection{Platoon Formation Drives Performance Gains}
The throughput, delay time, and energy consumption results all arise from the same root cause. The barrier controller converges the vehicles into a platoon, which is spaced in a manner that allows vehicles to cross the conflict point continuously, without mandatory stops. On the other hand, the FAS controller imposes idle time through its phase cycles. Despite its adaptive setup, the adaptive extension saturates beyond \(N=61\), resulting in the green time being extended to the maximum green time across all cycles, degrading the adaptive traffic light to a fixed-cycle traffic light. This difference explains why the throughput of the barrier controller scales near-linearly with vehicle count, while  the FAS controller saturates at approximately 0.32 veh/s beyond \(N=61\). 
\\\\
Delay time and energy consumption results follow from the continuous flow, which results in vehicles driving at target speed throughout the run, resulting in small deviations in lap time and the motor operating at low torque, in a high efficiency region. Conversely, the FAS controller does the opposite by imposing stop-and-go behavior. The stops result in decelerations, followed by rapid accelerations during the go, which drive the motor into operating at high-torque, in the low efficiency region. In addition, driving in queues forces the motor into driving at low torque and shaft speed, which is a second region of poor efficiency. Both effects have a greater impact at higher vehicle counts, since the traffic light already performs at maximum capacity, therefore the additional vehicles add directly to the queue. This explains why the gap in energy widens at \(N=71\) and \(N=81\), rather than scaling linearly. The energy reduction of maximum 44.81\% confirms that stop-and-go behavior is a primary driver of EV energy consumption at intersections.
\\\\
Safety is the only performance indicator showing opposite behavior. At low vehicle counts both controllers record zero violations. Whereas from \(N=71\), the barrier controller records safety violations during the phase before the continuous flow platoon is formed. This highlights a key limitation of the barrier controller setup: either the activation distance is too small, the cap of ±7 km/h is to small, or the barrier gain is set to gentle to separate vehicles with small inter-vehicle spacing. The FAS controller avoids this by having mutually exclusive phases which avoid simultaneous arrival. It is important to note, that violations are completely eliminated once steady state is reached, indicating that the barrier law itself is capable of safe operation. The violations reflect the initial spawn noise and the barrier corrections, rather than a fundamental barrier shortcoming.


\subsection{Simulation Limitations}
Although promising results have been found, this study also has limitations. First, the figure-eight track is a highly controlled environment, with a single conflict point, constant vehicle flow, no steering at the intersection, and absence of pedestrians or emergency vehicles. Therefore, performance at a multi-lane intersection with stochastic arrivals, pedestrians and emergency vehicles, and steering may differ. Second, the energy model relies on approximated efficiency maps from other literature, rather than manufacturer efficiency maps. The model is also subjected to severe assumptions, due to absence of data. Since both controllers are subjected to identical assumptions, the difference remains valid, but the produced outcomes should be interpreted with caution since they differ from actual values. Third, the barrier controller was only compared against a FAS controller based on Dutch standards according to \textcite{wilson2014handboek}. The book was published in 2014, so it is likely that more advanced controllers exist, which might close some of the performance gaps. Fourth, the simulation assumes that each vehicle is capable of perfectly estimating its own distance and the distance of the conflicting vehicle to the conflict point using onboard sensors. In practice, sensors are subjected to noise and estimation errors, which could reduce the accuracy of the distance measurement, which might affect performance and safety of the barrier controller.
